\medskip

\emph{Pour  chacune des affirmations suivantes, indiquer si elle est vraie ou fausse. Chaque
réponse doit être justifiée. Une réponse non justifiée ne rapporte aucun point :}

\medskip

\begin{enumerate}
\item \textbf{Affirmation :} La suite $u$ définie pour tout entier naturel $n$ par $u_n = \dfrac{(- 1)^n}{n + 1}$ est bornée.
\item \textbf{Affirmation :} Toute suite bornée est convergente.
\item \textbf{Affirmation :} Toute suite croissante tend vers $+\infty$.
\item Soit la fonction $f$ définie sur $\R$ par $f(x) = \ln \left(x^2 + 2x + 2\right)$.

\textbf{Affirmation :} La fonction $f$ est convexe sur l'intervalle $[-3~;~1]$.
\item On considère la fonction {\ttfamily \textbf{mystere}} définie ci-dessous qui prend une liste L de nombres en paramètre. 

On rappelle que \texttt{len}(L) renvoie la longueur, c'est-à-dire
le nombre d'éléments de la liste L.

\begin{center}
\begin{tabular}{|>{\ttfamily}l|}\hline
 def \textbf{mystere}(L) :\\
\qquad M = L[0]\\
\qquad \# \emph{On initialise M avec le premier élément de la liste }L\\
\qquad for i in range(len(L)) : \\
\qquad \quad if L[i] > M :\\
\qquad \qquad M = L[i]\\
\qquad return M \\ \hline
\end{tabular}
\end{center}

\textbf{Affirmation :} L'exécution de {\ttfamily \textbf{mystere}( [2, 3, 7,0 ,6, 3, 2, 0, 5])} renvoie 7.
\end{enumerate}
